% =====================================================================
%  3 Zusammenfassung der Testergebnisse
%  Gefundene Fehler, Behebung, offene Punkte (-> technische Schulden in
%  der Projektdokumentation).
%  Zahlen aus dem Lauf vom 27.11.2026 (junit.xml, pytest-cov, ruff).
% =====================================================================
\section{Zusammenfassung der Testergebnisse}\label{sec:testergebnisse}
Dieses Kapitel fasst die Ergebnisse des Testlaufs vom 27.11.2026 zusammen. Es zeigt die Ergebnisse je Teststufe, die Zeilenabdeckung und die Messwerte der nichtfunktionalen Anforderungen. Danach beschreibt es die Fehler, die Tests und Reviews während der Entwicklung gefunden haben, und die Punkte, die noch offen sind.

\subsection{Ergebnisübersicht}\label{sec:ergebnisuebersicht}
Tabelle~\ref{tab:testergebnisse} zeigt das Ergebnis je Teststufe. Die Spalte \enquote{Testfälle} zählt die Testfälle des Protokolls, die Spalte \enquote{pytest-Tests} die einzeln ausgeführten automatisierten Tests einschließlich aller Ausprägungen parametrisierter Tests. Die Laufzeit ist die Summe der Einzelzeiten aus dem Testbericht von pytest.

\begin{longtable}{@{}L{3.4cm}>{\raggedleft\arraybackslash}p{1.9cm}>{\raggedleft\arraybackslash}p{2.6cm}>{\raggedleft\arraybackslash}p{2.1cm}>{\raggedleft\arraybackslash}p{2.6cm}>{\raggedleft\arraybackslash}p{2.0cm}@{}}
  \caption{Übersicht der Testergebnisse}\label{tab:testergebnisse}\\
  \toprule
  \textbf{Teststufe} & \textbf{Testfälle} & \textbf{pytest-Tests} & \textbf{bestanden} & \textbf{fehlgeschlagen} & \textbf{Laufzeit} \\
  \midrule
  \endfirsthead
  \toprule
  \textbf{Teststufe} & \textbf{Testfälle} & \textbf{pytest-Tests} & \textbf{bestanden} & \textbf{fehlgeschlagen} & \textbf{Laufzeit} \\
  \midrule
  \endhead
  \bottomrule
  \endfoot
  Unit-Tests & 16 & 141 & 141 & 0 & 0,2~s \\
  Integrationstests & 15 & 176 & 176 & 0 & 5,8~s \\
  \quad davon Oberfläche & -- & 93 & 93 & 0 & 2,7~s \\
  \quad davon Lasttests & 1 & 4 & 4 & 0 & 2,7~s \\
  Systemtests & 10 & 9 + 1 manuell & 9 + 1 & 0 & 18,0~s \\
  \midrule
  Summe & 41 & 326 + 1 manuell & 326 + 1 & 0 & 24,0~s \\
\end{longtable}
\quelle{Eigene Darstellung.}

Alle 326 automatisierten Tests haben bestanden. Kein Test wurde übersprungen, und es gab keine Fehler beim Aufbau der Tests. Der manuelle Testfall TC-S-10 hat ebenfalls bestanden, mit der in Abschnitt~\ref{sec:fehler} genannten Einschränkung. Der gesamte Lauf mit Abdeckungsmessung dauerte 25,6~s. Gut ein Drittel davon entfällt auf TC-S-07, der zehn Python-Prozesse nacheinander startet und hart beendet (8,8~s). Die Unit-Tests brauchen zusammen nur etwa 0,2~s und lassen sich deshalb nach jeder Änderung ausführen.

Die Zahl der Tests je Stufe spiegelt die Architektur wider. Die meisten Unit-Tests betreffen die Regeln (33) und die Gültigkeitsregeln (48 einschließlich Parametrisierung), weil dort die meisten Fallunterscheidungen der Spezifikation liegen. Die Integrationstests verteilen sich auf alle Anwendungsservices und alle Seiten der Oberfläche. Die neun automatisierten Systemtests decken die Abläufe ab, die die Prüfkriterien der nichtfunktionalen Anforderungen verlangen, und zusätzlich die beiden wichtigsten Nutzungsabläufe: den ersten Start mit Demodaten und das Erfassen einer Aktivität. 53 der 326 Tests kamen erst nach der Abschlussprüfung der Software hinzu (TC-U-16, TC-I-12 bis TC-I-15 und ein Test in TC-I-08). 44 davon sind Oberflächentests, weil die meisten gefundenen Fehler die Oberfläche betrafen.

\subsection{Abdeckung und Messwerte}\label{sec:abdeckung}
Tabelle~\ref{tab:abdeckung} zeigt die Zeilenabdeckung, die pytest-cov im vollständigen Lauf gemessen hat. Für NFA07 zählt die Abdeckung der Bewertungs- und Regellogik durch die Unit-Tests allein. Sie wurde deshalb zusätzlich in einem Lauf nur mit \texttt{tests/unit} gemessen.

\begin{longtable}{@{}L{7.3cm}>{\raggedleft\arraybackslash}p{2.6cm}>{\raggedleft\arraybackslash}p{3.2cm}>{\raggedleft\arraybackslash}p{2.6cm}@{}}
  \caption{Zeilenabdeckung des Quellcodes}\label{tab:abdeckung}\\
  \toprule
  \textbf{Bereich} & \textbf{Anweisungen} & \textbf{nicht ausgeführt} & \textbf{Abdeckung} \\
  \midrule
  \endfirsthead
  \toprule
  \textbf{Bereich} & \textbf{Anweisungen} & \textbf{nicht ausgeführt} & \textbf{Abdeckung} \\
  \midrule
  \endhead
  \bottomrule
  \endfoot
  \multicolumn{4}{@{}l}{\textit{Alle Tests (326)}} \\*
  \texttt{balanceos} gesamt & 4.157 & 173 & 96~\% \\
  \quad Domäne (\texttt{domaene}) & 621 & 6 & 99~\% \\
  \quad Anwendung (\texttt{anwendung}) & 544 & 8 & 99~\% \\
  \quad Persistenz (\texttt{persistenz}) & 257 & 4 & 98~\% \\
  \quad Präsentation (\texttt{praesentation}) & 2.730 & 151 & 94~\% \\
  \midrule
  \multicolumn{4}{@{}l}{\textit{Nur Unit-Tests (141), Prüfkriterium NFA07}} \\*
  \texttt{domaene/bewertung.py} & 41 & 0 & 100~\% \\
  \texttt{domaene/regeln.py} & 183 & 4 & 98~\% \\
  Domäne gesamt & 621 & 7 & 99~\% \\
\end{longtable}
\quelle{Eigene Darstellung.}

Die Bewertungs- und Regellogik liegt mit 100~\% und 98~\% deutlich über dem Grenzwert von 80~\%. Die niedrigste Abdeckung hat die Präsentationsschicht. Die Startfunktion in \texttt{praesentation/app.py} (81~\%) und \texttt{\_\_main\_\_.py} (0~\%) werden zwar in TC-S-08 ausgeführt, aber in einem eigenen Prozess, den pytest-cov nicht misst. Die übrigen Lücken liegen vor allem im Wochenkalender (89~\%) bei Doppelklick, Tooltip und dem Zeichnen ausgefallener Termine sowie in Dashboard und Erfassungsdialog (je 92~\%).

Für NFA02 verlangt das Prüfkriterium, dass alle Zeilen der Tabelle \enquote{Gültigkeitsregeln für Eingabefelder} des Spezifikationsdokuments durch bestandene Testfälle abgedeckt sind. Tabelle~\ref{tab:validierungsabdeckung} ordnet die 26 Zeilen den Testfällen zu.

\begin{longtable}{@{}L{3.5cm}L{4.3cm}L{3.5cm}L{4.3cm}@{}}
  \caption{Abdeckung der Gültigkeitsregeln für Eingabefelder}\label{tab:validierungsabdeckung}\\
  \toprule
  \textbf{Feld} & \textbf{Testfälle} & \textbf{Feld} & \textbf{Testfälle} \\
  \midrule
  \endfirsthead
  \toprule
  \textbf{Feld} & \textbf{Testfälle} & \textbf{Feld} & \textbf{Testfälle} \\
  \midrule
  \endhead
  \bottomrule
  \endfoot
  Titel & TC-U-09, TC-I-03, TC-S-03 & Name & TC-U-10, TC-I-02 \\
  Lebensbereich & TC-U-09, TC-I-03 & Typ, Farbe & TC-I-02 \\
  Datum & TC-I-03, TC-S-02 & Aktiv & TC-U-10, TC-I-02, TC-I-14 \\
  Beginn, Ende & TC-U-04, TC-U-09, TC-I-03 & Soll-Stunden & TC-U-10, TC-I-14 \\
  Dauer & TC-U-09, TC-S-03 & Summe Soll & TC-U-10, TC-I-02, TC-I-14 \\
  Status & TC-U-09, TC-I-03 & Gewicht & TC-U-10, TC-I-02 \\
  Zufriedenheit, Belastung, Erholung & TC-U-09, TC-I-03 & Gewichtung $\alpha$ & TC-U-11, TC-I-02 \\
  Bewegungsart & TC-U-09, TC-I-03 & Datum WHO-5 & TC-U-11, TC-I-10 \\
  Notiz & TC-U-09 & Aussagen 1--5 & TC-U-11, TC-I-10, TC-I-14 \\
  Wochentage & TC-U-10, TC-I-04 & Leitfragen, Stimmung & TC-U-11, TC-I-10, TC-I-14 \\
  Enddatum & TC-U-10, TC-I-04 & Format & TC-I-08, TC-S-04 \\
  Zeitraum & TC-U-11, TC-I-08 & Speicherort & TC-I-08 \\
  Demodaten & TC-I-09 & Bestätigung & TC-U-11, TC-I-09, TC-I-14, TC-S-05 \\
\end{longtable}
\quelle{Eigene Darstellung.}

Alle 26 Zeilen werden mit ungültigen Eingaben oder Grenzwerten geprüft. Bei der Zeile \enquote{Datum} prüft TC-I-03, dass das Kalenderfeld das Format TT.MM.JJJJ verwendet und ein ungültiges Datum wie den 31.02.2026 nicht übernimmt.

Tabelle~\ref{tab:nfa-messwerte} stellt für jede nichtfunktionale Anforderung den Grenzwert aus dem Anforderungsdokument dem gemessenen Wert gegenüber.

\begin{longtable}{@{}L{1.4cm}L{4.8cm}L{4.9cm}L{2.0cm}L{2.2cm}@{}}
  \caption{Messwerte der nichtfunktionalen Anforderungen}\label{tab:nfa-messwerte}\\
  \toprule
  \textbf{NFA} & \textbf{Grenzwert} & \textbf{Messwert am 27.11.2026} & \textbf{Testfall} & \textbf{Ergebnis} \\
  \midrule
  \endfirsthead
  \toprule
  \textbf{NFA} & \textbf{Grenzwert} & \textbf{Messwert am 27.11.2026} & \textbf{Testfall} & \textbf{Ergebnis} \\
  \midrule
  \endhead
  \bottomrule
  \endfoot
  NFA01 & kein englischer Standardtext; 100~\% der Eingabefelder beschriftet; Erfassung mit 1~Klick & 0 Verstöße; 7 von 7 Seiten mit 1~Klick & TC-S-09 & erfüllt \\
  NFA02 & 100~\% der Zeilen der Validierungstabelle & 26 von 26 Zeilen & TC-U-09 bis TC-U-11, TC-I-03, TC-I-14, TC-S-03 & erfüllt \\
  NFA03 & Start höchstens 5~s; Dashboard bei 3.000 Aktivitäten höchstens 1~s & Start 0,82~s; Dashboard 5~ms & TC-S-08, TC-I-11 & erfüllt \\
  NFA04 & 0 Verluste in 10 Versuchen & 0 von 10 verloren; Integritätsprüfung \enquote{ok} & TC-S-07 & erfüllt \\
  NFA05 & 0 Netzwerkzugriffe; 0 Nutzdaten nach dem Löschen & 0 Zugriffe; 0 Datensätze in 6 Tabellen & TC-S-06, TC-S-05 & erfüllt \\
  NFA06 & 100~\% der Rechenbeispiele & 3 von 3 & TC-U-01 & erfüllt \\
  NFA07 & Zeilenabdeckung mindestens 80~\% & \texttt{bewertung.py} 100~\%, \texttt{regeln.py} 98~\% & TC-U-14 & erfüllt \\
  NFA08 & 0 Befunde; 100~\% Docstrings & 0 Befunde; 312 von 312 & TC-U-15 & erfüllt \\
  NFA09 & Hauptfenster nach höchstens 10~s ohne installiertes Python & 3,8~s, 3,3~s und 3,2~s auf einem Rechner mit Python & TC-S-10 & vorläufig erfüllt \\
  NFA10 & UTF-8 mit BOM, Semikolon, 16 Spalten je Zeile; JSON gültig & alle Bedingungen erfüllt & TC-S-04 & erfüllt \\
  NFA11 & 10 von 10 Regeln vollständig begründet & 10 von 10 & TC-U-08 & erfüllt \\
  NFA12 & Dashboard und Statistik je höchstens 2~s; Datei unter 50~MB & Dashboard 5~ms; Statistik 10~ms; 3,2~MB & TC-I-11 & erfüllt \\
\end{longtable}
\quelle{Eigene Darstellung.}

Die Leistungswerte liegen weit unter den Grenzwerten. Die Lasttests messen die Neuberechnung im Auswertungsservice, also Wochenbewertung, Hinweise und Verlauf. Das Zeichnen der Oberfläche ist darin nicht enthalten. Die Werte stammen aus einem einzelnen Lauf und schwanken von Lauf zu Lauf um einige Millisekunden. Die Dauer hängt kaum von der Gesamtzahl der Aktivitäten ab, weil der Auswertungsservice aus der Datenbank nur die Aktivitäten der benötigten Wochen liest. Die Startzeit von 0,82~s gilt für den Start mit \texttt{python -m balanceos}. Die ausführbare Datei braucht etwa 3 bis 4~s, weil sie ihre Laufzeitumgebung bei jedem Start in einen temporären Ordner entpackt.

\subsection{Gefundene Fehler und offene Punkte}\label{sec:fehler}
Im abschließenden Lauf ist kein Test fehlgeschlagen. Während der Entwicklung haben Tests und Reviews jedoch mehrere Fehler gefunden. Tabelle~\ref{tab:fehler} beschreibt sie mit Ursache und Behebung. Zu jedem behobenen Fehler gibt es einen Regressionstest, der sein erneutes Auftreten erkennt. Die Angaben lassen sich in der Versionsgeschichte des Repositorys nachvollziehen (Commits \texttt{7bf7a31} vom 23.11.2026 sowie \texttt{69ca627}, \texttt{7ccd33a} und \texttt{64d7188} vom 27.11.2026).

Die Fehler F6 bis F13 fand die Abschlussprüfung der Software in Phase~3, nachdem alle automatisierten Tests bereits bestanden hatten. Dabei wurde die laufende Anwendung in Bedienabläufen über mehrere Schritte geprüft, zum Beispiel einen Serientermin verschieben und die Änderung danach auf alle zukünftigen Termine anwenden. Solche Abfolgen deckten die automatisierten Tests bis dahin nicht ab. Jeder Fehler wurde zuerst als fehlschlagender Test nachgestellt und erst danach behoben. F14 kam nach der Abschlussprüfung hinzu.

\begin{longtable}{@{}L{0.7cm}L{4.4cm}L{4.6cm}L{6.0cm}@{}}
  \caption{Gefundene und behobene Fehler}\label{tab:fehler}\\
  \toprule
  \textbf{Nr.} & \textbf{Fehler und Fundstelle} & \textbf{Ursache} & \textbf{Behebung und Regressionstest} \\
  \midrule
  \endfirsthead
  \toprule
  \textbf{Nr.} & \textbf{Fehler und Fundstelle} & \textbf{Ursache} & \textbf{Behebung und Regressionstest} \\
  \midrule
  \endhead
  \bottomrule
  \endfoot
  F1 & Ein gelöschter oder auf einen anderen Tag verschobener Serientermin wurde bei der nächsten Änderung der Serie neu erzeugt. Gefunden im Abschlussreview der Phase~2. & Löschen und Verschieben eines einzelnen Termins legten keine Ausnahme an. Die Serienänderung kannte das ursprüngliche Datum deshalb nicht als bearbeitet. & Löschen und Verschieben tragen das ursprüngliche Datum als Ausnahme der Serie ein (GR18), sodass eine Serienänderung es auslässt (GR19). Verschobene Termine lassen sich überspringen. Tests in TC-I-04, u.\,a. \pyt{test_serie_aendern_erzeugt_geloeschten_termin_nicht_neu}. \\
  \midrule
  F2 & R05 zählte Aktivitäten deaktivierter Lebensbereiche mit. Gefunden im Abschlussreview der Phase~2. & Die Hilfsfunktion, die die durchgeführten Aktivitäten einer Woche liefert, filterte nicht nach aktiven Bereichen. & Filter auf aktive Bereiche ergänzt (GR14). Test \pyt{test_r05_ignoriert_deaktivierten_bereich} (TC-U-05). \\
  \midrule
  F3 & R07 meldete Arbeit über dem Soll, obwohl für Arbeit kein Soll festgelegt war. Gefunden im Abschlussreview der Phase~2. & Bei Soll 0~h war die Bedingung \enquote{Ist über Soll} schon mit jeder Arbeitsstunde erfüllt. Soll 0~h bedeutet aber \enquote{kein Soll festgelegt} (GR22). & R07 entfällt bei einem Soll von 0~h. Test \pyt{test_r07_ohne_soll_kein_hinweis} (TC-U-06). \\
  \midrule
  F4 & Eingabefelder ohne verknüpfte Beschriftung: Filter der Seite \enquote{Aktivitäten}, Felder des Wunschprofils, Enddatum der Serie. Gefunden durch TC-S-09. & Die Beschriftung stand als Text neben dem Feld, war aber nicht mit ihm verknüpft. & Beschriftungen mit dem Feld verknüpft bzw. einen zugänglichen Namen gesetzt. TC-S-09 prüft alle sichtbaren Felder auf allen Seiten. \\
  \midrule
  F5 & Rückfragen von Qt zeigten die englischen Knöpfe \enquote{Yes} und \enquote{No}. Gefunden bei der Arbeit an den Systemtests in Phase~3. & Die Qt-Übersetzung der Standardtexte (\texttt{qtbase\_de}) wurde beim Start nicht geladen. & Die Übersetzung wird beim Start geladen und in die ausführbare Datei gepackt; fehlt sie, startet die Anwendung trotzdem. Tests in \pyt{tests/gui/test_uebersetzung.py} (TC-I-05). \\
  \midrule
  F6 & Wurde ein Serientermin im Dialog auf einen anderen Tag verschoben und dabei \enquote{auf alle zukünftigen Termine anwenden} gewählt, entstand am alten Tag ein zweiter Termin. & Der Dialog speicherte die Serie in dem Stand, den er beim Öffnen geladen hatte. Die Ausnahme, die das Verschieben kurz zuvor angelegt hatte, ging dabei verloren. & Der Dialog lädt die Serie vor dem Speichern neu, und die Serienänderung vereinigt die gespeicherten Ausnahmen mit den übergebenen (GR18, GR19). Tests in TC-I-12. \\
  \midrule
  F7 & Ungespeicherter Text der Wochenreflexion ging verloren, wenn der WHO-5-Fragebogen ausgewertet oder die Woche gewechselt wurde. & Die Seite lud bei jeder Aktualisierung die gespeicherte Reflexion neu und überschrieb damit die Eingabefelder. & Die Auswertung lässt ungespeicherten Text stehen. Vor einem Wochenwechsel fragt eine Rückfrage, ob gespeichert, verworfen oder abgebrochen werden soll. Tests \pyt{test_a2_*} in TC-I-14. \\
  \midrule
  F8 & \enquote{Termin überspringen} auf der Seite \enquote{Aktivitäten} löschte ohne Rückfrage, auch durchgeführte Termine mit Bewertung und Notiz. & Der Knopf war für jeden Termin einer Serie aktiv und führte das Überspringen sofort aus. & Überspringen ist nur für geplante Serientermine möglich und erst nach einer Rückfrage (US08). Tests in TC-I-12. \\
  \midrule
  F9 & Namen und Titel mit \enquote{<} wurden als HTML dargestellt, z.\,B. erschien \enquote{<b>Hobbys</b>} fett und ohne Tags. & Beschriftungen, Rückfragen und Tooltips erkennen HTML automatisch, die Texte wurden aber ungeschützt übergeben. & Texte werden als Klartext gesetzt oder vorher maskiert. Tests \pyt{test_b1_*} und \pyt{tests/gui/test_review_phase3_klartext.py} in TC-I-15. \\
  \midrule
  F10 & Das Hauptfenster passte nicht auf kleine Bildschirme. & Die Mindestgröße war fest auf 1.280 $\times$ 720 Pixel gesetzt. & Mindestgröße 960 $\times$ 520 Pixel; die Seiten liegen in einem Bildlaufbereich, die Startgröße richtet sich nach der Bildschirmfläche. Tests \pyt{test_mindestgroesse} (TC-I-05) und \pyt{test_b2_*} (TC-I-15). \\
  \midrule
  F11 & Ein Datenbankfehler im laufenden Betrieb, etwa eine gesperrte Datei, blieb ohne Meldung. & Unbehandelte Fehler in der Oberfläche wurden nur auf der Konsole ausgegeben, die die ausführbare Datei nicht hat. & Eine zentrale Fehlerbehandlung zeigt eine deutsche Meldung und schreibt die Einzelheiten in \texttt{fehler.log}. Tests \pyt{test_b3_*} in TC-I-13. \\
  \midrule
  F12 & Datenbankdatei: Nach \enquote{Alle Daten löschen} standen gelöschte Inhalte bis zum Programmende im Write-Ahead-Log; bei einem Fehler beim Öffnen blieb die Verbindung offen; die Anwendung ließ sich zweimal mit denselben Daten starten. & Das Log wurde erst beim Beenden geleert, die Verbindung im Fehlerfall nicht geschlossen, und eine Sperre gegen eine zweite Instanz fehlte. & Nach dem Löschen wird das Log sofort geleert (NFA05, GR29), die Verbindung bei einem Fehler geschlossen und der Datenordner gesperrt. Tests in TC-I-13. \\
  \midrule
  F13 & Kleinere Fehler in Texten und Darstellung: 6,98~h erschien in R03 als \enquote{7 h}; Einzahl fehlte (\enquote{1 Aktivitäten}, \enquote{Noch offen: 1 Aussagen.}); das Dashboard zeigte einen Strich statt \enquote{zu wenig Daten} und keine Abweichung Ist $-$ Soll; Meldungen am Feld fehlten für Leitfragen, Löschbestätigung und Lebensbereich; die Wochenansicht zeigte den Überlauf vom Sonntag der Vorwoche nicht. & Verschiedene Einzelursachen in Regeltexten, Seiten und Wochenkalender. & Einzeln behoben. Tests in TC-U-16, TC-I-14 und TC-I-15; angepasst wurden \pyt{test_mindestgroesse}, \pyt{test_unzureichende_daten_zeigen_zu_wenig_daten} und TC-S-02. \\
  \midrule
  F14 & Im CSV-Export konnte ein Titel wie \enquote{=SUMME(1;2)} beim Öffnen in einer Tabellenkalkulation als Formel ausgeführt werden. & Textfelder wurden unverändert geschrieben. & Bereichsname, Titel und Notiz erhalten ein vorangestelltes Apostroph, wenn sie mit =, +, -, @, Tabulator oder Wagenrücklauf beginnen. Test \pyt{test_csv_textfelder_mit_formelzeichen_erhalten_apostroph} (TC-I-08). \\
\end{longtable}
\quelle{Eigene Darstellung.}

Ein weiteres Problem betraf das Testwerkzeug selbst. Die Funktion \texttt{QTest.keyClicks} von Qt~6.12 bricht den Prozess ab, wenn sie ein Zeichen ohne eigene Qt-Taste eingeben soll, etwa \enquote{ü} oder das deutsche Anführungszeichen. Die Systemtests senden solche Zeichen deshalb als Tastenereignis mit Text an das Feld, so wie es eine Tastatur mit deutschem Layout erzeugt (Hilfsfunktion \texttt{zeichen\_eingeben}). Die Anwendung selbst war davon nicht betroffen.

Die Aussagekraft der Tests hat Grenzen. TC-S-10 lief nur auf dem Entwicklungsrechner, auf dem Python installiert ist. Die ausführbare Datei nutzt dieses Python nicht, der Nachweis auf einem Rechner ohne Python fehlt aber. NFA09 gilt deshalb nur als vorläufig erfüllt. Einen Usability-Test mit Personen aus der Zielgruppe gab es nicht, und die Oberflächentests laufen ohne sichtbaren Bildschirm. Außerdem zeigen die Tests nur, dass Scores und Hinweise der Spezifikation entsprechen, nicht, dass der Score die Balance einer Person treffend abbildet. Die Projektdokumentation führt diese Punkte im Kapitel \enquote{Anforderungserfüllung und technische Schulden} als TD02 bis TD04 und bewertet sie im Fazit.
